\pdfbookmark[1]{Explanations of the plots and numerical readings for the scheduled observations. 1837.}{erlaeuterungen-termine-1837}
\begin{center}\large\textit{Explanations of the plots and numerical readings for the scheduled observations.}\par
[\textit{An excerpt.}]\end{center}
\begin{center}\small Results from the Observations of the Magnetic Association.\quad 1837.\quad VIII.\end{center}

Seven twenty-four-hour sessions were held in 1837, since an extraordinary one on 31 August was added to the six regular ones. The numbers given in the ---\quad---\quad--- tables contain 80 series of observations of declination changes from 16 different observing sites, namely Altona, Augsburg, Berlin, Breda, Breslau, Copenhagen, Dublin, Freiberg, Göttingen, Leipzig, Milan, Marburg, Munich, St Petersburg, Stockholm and Uppsala. We also received some other series of observations that arrived too late to be printed with the rest.
\begin{center}\small
---\enspace---\enspace---\enspace---\enspace---\enspace---\enspace---\enspace---\enspace---\enspace---\enspace---\enspace---\enspace---\enspace---\enspace---
\end{center}
The extraordinary session on 31 August was arranged on receiving news that Professor \textsc{Parrot} would be observing variations in magnetic declination at the North Cape that day (and on some preceding days). An invitation to take part was therefore circulated as far as the short notice allowed. This brought in some quite interesting series of observations, but we have yet to learn of the observations from the North Cape itself. ---\quad---\quad---\quad---\quad---\quad---\quad---

For the November session the usual arrangement was changed and the session moved to the 13th. This followed a conversation with Mr.\ von \textsc{Humboldt} about the possibility that unusual magnetic movements might also occur on the days of the month that in several earlier years had been marked by an extraordinary number of shooting
\SourceLeaf{583}stars. This expectation was not borne out, however, in that the magnetic movements during those twenty-four hours, though very considerable, were no greater than in many earlier sessions at any other time of year. On the preceding and the following evening, by contrast, very strong and rapidly changing anomalies in magnetic declination were observed at several sites. But we have no grounds for assuming a connection between those anomalies and the shooting stars, since the anomalies are simply the usual accompaniments of the northern lights, and very brilliant northern lights did indeed occur on both those nights\expandafter\def\csname HistoricalNoteText583\endcsname{We have received the magnetic observations for 12 November from Uppsala, Leipzig, Breslau and Milan, and for 14 November from Uppsala, Dublin, Berlin, Breslau and Milan. At some points the resemblance between the movements is unmistakable, while in others it is only faintly discernible. Here, however, a point already noted on another occasion bears repeating. Under such circumstances the movements change far too rapidly for observations taken at five-minute intervals, let alone at still longer intervals, to give a faithful picture of them.}\HistoricalNote{583}.
\begin{center}\small
---\enspace---\enspace---\enspace---\enspace---\enspace---\enspace---\enspace---\enspace---\enspace---\enspace---\enspace---\enspace---\enspace---\enspace---\enspace---
\end{center}
In the July, August and November sessions the variations in intensity at Göttingen were now also observed in full with the bifilar magnetometer. The tables do not give the directly observed scale readings themselves. but their differences from the largest value that occurred in each session. Since the transverse position used in the first two sessions was the one in which increasing scale readings correspond to decreasing intensities, the numbers here indicate by how much each intensity exceeded the smallest intensity of the session, measured in units of which the smallest intensity itself amounted to 22000 in the July session. For the time being, as long as such observations are made at only one site, there is no need to give the \textit{absolute} value of the scale divisions with the utmost precision. For the August session, therefore, no new determinations were made for that purpose. They were made before the November session, however. The change in the absolute number is connected with the loss of magnetism the bar suffered in the four months. It must be noted, though, that in the November session the numbers mean
\SourceLeaf{584}how much smaller each intensity is than the largest, the latter itself being taken as $=18290$. For in that session the bar was in the opposite transverse position, in which intensity and scale readings increase together. To give the numbers a consistent meaning, the individual scale readings observed directly would therefore have had to be compared with the minimum rather than with the maximum. This was overlooked by mistake.

The observed changes in intensity are plotted in the plates ---\quad---\quad---\quad---. They are shown first in one curve, below which the curve for the simultaneous changes in declination at Göttingen is repeated. This makes the observation made above [p.\ 365 of this volume] evident to the eye, namely that when there are strong disturbances of declination there are usually also strong anomalies in intensity. Second, the course of the changes in both elements is also combined in one curve, which gives a visual picture of the changes in the horizontal component of the Earth’s magnetic force during each session. To avoid confusion from the many crossings, however, the movements of the July and August sessions had to be drawn in two pieces and those of the November session in three. For greater ease of inspection, one half of each piece is drawn with solid lines and the other half with dots. After what was already said above [p.\ 366 of this volume], these plots will require no further explanation.

Only a few further remarks on the results themselves can find a place here. The extraordinarily close resemblance of the simultaneous declination movements at different sites, usually down to the finest shades of difference, is confirmed again here just as beautifully as in the previous year’s observations. Yet considerable differences can now also be discerned here and there, especially in the sessions in which the observations cover a still wider area, although that area still seems too small and the number of sites far apart too few to warrant any conclusions yet about where the causes of the individual movements are located. A closer examination of some individual movements, especially from sessions in which changes in intensity were observed at Göttingen at the same time, could indeed give rise
\SourceLeaf{585}to all sorts of remarks and even general reflections. We would, however, leave these to our readers and add this reminder instead. Where any dissimilarities appear, the external circumstances must first be considered carefully before the dissimilarities are made the basis of daring conjectures. One example is the small rise noticed in the plots of the August session ---\quad---\quad---\quad--- at $18^{\mathrm u}5'$ at most of the observing sites, most strongly at the northernmost, Uppsala. Its absence at Dublin, or the fact that only a faint trace of it can be seen there, is indeed remarkable, since there is no reason to doubt the correctness of the observation itself. It would give us occasion for interesting reflections, especially when considered together with its full manifestation at Göttingen ---\quad---, if it were at all appropriate to go into such matters already here. But the fact that this rise is also absent at Copenhagen has absolutely no significance, because observations in Copenhagen were made only every 10 minutes in this session, so there was no observation at all at $18^{\mathrm u}5'$.

We refrain here from commenting on the labyrinthine forms that the magnetic observations present when the movements of declination and intensity are combined in one curve ---\quad---\quad---\quad---\quad---, solely because there is good reason to hope that much richer material will soon be available. Anyone who would like to consider them independently in the meantime need at least feel no hesitation on account of doubts about the reality of the movements of intensity indicated by the bifilar magnetometer. Indeed, such doubts have become quite untenable. As early as the March session of the present year, 1838, the simultaneous movements of intensity were observed with similar bifilar instruments at \textit{three other sites} as well as in Göttingen, and they showed an agreement just as admirable as the one we have been accustomed to finding for four years in the movements of declination. Further details, however, will have to wait for the publication of the results of the observations of 1838.

